% p341_1; source PDF page 341; manually verified editable table.
\begin{tikzpicture}[x=1cm,y=-1cm,inner sep=0pt,outer sep=0pt]
\path[use as bounding box] (0,0) rectangle (16.000,8.725);
\draw[line width=.4pt] (0.000,0.000) rectangle (8.350,0.865);
\node[align=center,text width=8.090cm,font={\linespread{1}\fontsize{10}{13.50}\selectfont\bfseries}] at (4.175,0.432) {同一对象};
\draw[line width=.4pt] (8.350,0.000) rectangle (16.000,0.865);
\node[align=center,text width=7.390cm,font={\linespread{1}\fontsize{10}{13.50}\selectfont\bfseries}] at (12.175,0.432) {不同对象};
\draw[line width=.4pt] (0.000,0.865) rectangle (8.350,3.530);
\node[align=left,text width=8.090cm,font={\linespread{1}\fontsize{10}{13.50}\selectfont}] at (4.175,2.198) {伪造+出售、购买、运输、持有、使用=伪造货币罪\\{}（这是“先干什么、后干什么”的串联关系，前罪吸收后罪）};
\draw[line width=.4pt] (8.350,0.865) rectangle (16.000,3.530);
\node[align=left,text width=7.390cm,font={\linespread{1}\fontsize{10}{13.50}\selectfont}] at (12.175,2.198) {伪造A假币+出售、购买、运输、持有、使用B假币，并罚\\{}（前后两个罪是独立罪名，所以并罚）（2013年第14题）};
\draw[line width=.4pt] (0.000,3.530) rectangle (8.350,5.645);
\node[align=left,text width=8.090cm,font={\linespread{1}\fontsize{10}{13.50}\selectfont}] at (4.175,4.588) {出售+购买+运输，定出售、购买、运输假币罪};
\draw[line width=.4pt] (8.350,3.530) rectangle (16.000,5.645);
\node[align=left,text width=7.390cm,font={\linespread{1}\fontsize{10}{13.50}\selectfont}] at (12.175,4.588) {出售A假币+购买B假币+运输C假币，定出售、购买、运输假币罪，不并罚（这是选择性罪名，无论对象是否同一，都不并罚）};
\draw[line width=.4pt] (0.000,5.645) rectangle (8.350,7.210);
\node[align=left,text width=8.090cm,font={\linespread{1}\fontsize{10}{13.50}\selectfont}] at (4.175,6.428) {出售/运输+使用，并罚（这是“边干什么、边干什么”的并联关系，所以并罚）};
\draw[line width=.4pt] (8.350,5.645) rectangle (16.000,7.210);
\node[align=left,text width=7.390cm,font={\linespread{1}\fontsize{10}{13.50}\selectfont}] at (12.175,6.428) {出售/运输A假币+使用B假币，并罚（前后两个罪是独立罪名，所以并罚）（2010年第13题）};
\draw[line width=.4pt] (0.000,7.210) rectangle (8.350,8.725);
\node[align=left,text width=8.090cm,font={\linespread{1}\fontsize{10}{13.50}\selectfont}] at (4.175,7.968) {购买+使用=购买假币罪（这是串联关系，前罪吸收后罪)};
\draw[line width=.4pt] (8.350,7.210) rectangle (16.000,8.725);
\node[align=left,text width=7.390cm,font={\linespread{1}\fontsize{10}{13.50}\selectfont}] at (12.175,7.968) {购买A假币+使用B假币，并罚（前后两个罪是独立罪名，所以并罚）};
\end{tikzpicture}
